% Part: first-order-logic
% Chapter: sequent-calculus
% Section: derivations

\documentclass[../../../include/open-logic-section]{subfiles}

\begin{document}

\iftag{FOL}
      {\olfileid[id]{fol}{seq}{der}}
      {\olfileid[id]{pl}{seq}{der}}

\olsection{\usetoken{P}{derivation}}

\begin{explain}
Kita telah menjelaskan bentuk sekuen awal dan memberikan aturan-aturan
inferensi. !!^{derivation}s dalam kalkulus sekuen dibangkitkan secara
induktif dari keduanya: setiap !!{derivation} merupakan sekuen awal itu
sendiri, atau terdiri atas satu atau dua !!{derivation}s yang diikuti oleh
suatu inferensi.
\end{explain}

\begin{defn}[!!^{derivation} $\Log{LK}$]
Sebuah \emph{$\Log{LK}$-!!{derivation}} dari sekuen~$S$ adalah pohon
berhingga yang terdiri atas sekuen-sekuen dan memenuhi syarat berikut:
\begin{enumerate}
\item Sekuen-sekuen paling atas pada pohon tersebut merupakan sekuen awal.
\item Sekuen paling bawah pada pohon tersebut adalah~$S$.
\item Setiap sekuen pada pohon tersebut selain $S$ merupakan premis dari
  penerapan aturan inferensi yang benar, dan kesimpulan aturan itu berada
  tepat di bawah sekuen tersebut pada pohon.
\end{enumerate}
Kita lalu mengatakan bahwa $S$ adalah \emph{sekuen akhir} dari
!!{derivation} tersebut dan bahwa $S$ \emph{!!{derivable} dalam
$\Log{LK}$} (atau $\Log{LK}$-!!{derivable}).
\end{defn}

\begin{ex}
Setiap sekuen awal, misalnya $!C \Sequent !C$, merupakan
!!a{derivation}. Kita dapat memperoleh !!a{derivation} baru darinya
dengan menerapkan, misalnya, aturan $\LeftR{\Weakening}$,
\begin{prooftree}
\Axiom$ \Gamma \fCenter \Delta $
\RightLabel{\LeftR{\Weakening}}
\UnaryInf$ !A, \Gamma \fCenter \Delta$
\end{prooftree}
Namun, aturan tersebut dimaksudkan untuk berlaku umum: kita dapat
mengganti $!A$ dalam aturan itu dengan sebarang !!{sentence}, misalnya
juga dengan~$!D$. Jika premisnya dicocokkan dengan sekuen awal kita
$!C \Sequent !C$, ini berarti $\Gamma$ dan $\Delta$ masing-masing
hanyalah~$!C$, sehingga kesimpulannya adalah $!D, !C \Sequent !C$.
Jadi, berikut ini merupakan !!a{derivation}:
\begin{prooftree}
\Axiom$ !C \fCenter !C $
\RightLabel{\LeftR{\Weakening}}
\UnaryInf$ !D, !C \fCenter !C$
\end{prooftree}
Sekarang kita dapat menerapkan aturan lain, misalnya
$\LeftR{\Exchange}$, yang memungkinkan kita menukar dua !!{sentence}s
di sebelah kiri. Jadi, berikut ini juga merupakan !!a{derivation} yang
benar:
\begin{prooftree}
\Axiom$ !C \fCenter !C $
\RightLabel{\LeftR{\Weakening}}
\UnaryInf$ !D, !C \fCenter !C$
\RightLabel{\LeftR{\Exchange}}
\UnaryInf$ !C, !D \fCenter !C$
\end{prooftree}
Dalam penerapan aturan ini, yang diberikan sebagai
\begin{prooftree}
\Axiom$ \Gamma, !A, !B, \Pi \fCenter \Delta $
\RightLabel{\LeftR{\Exchange}}
\UnaryInf$ \Gamma, !B, !A, \Pi \fCenter \Delta,$
\end{prooftree}
$\Gamma$ dan $\Pi$ sama-sama kosong, $\Delta$ adalah $!C$, dan peran
$!A$ serta $!B$ masing-masing dimainkan oleh $!D$ dan~$!C$. Dengan cara
yang hampir sama, kita juga melihat bahwa
\begin{prooftree}
\Axiom$ !D \fCenter !D $
\RightLabel{\LeftR{\Weakening}}
\UnaryInf$ !C, !D \fCenter !D$
\end{prooftree}
merupakan !!a{derivation}. Sekarang kita dapat mengambil kedua
!!{derivation}s ini dan menggabungkannya dengan $\RightR{\land}$. Aturan
tersebut adalah
\begin{prooftree}
\Axiom$\Gamma \fCenter \Delta, !A$
\Axiom$ \Gamma \fCenter \Delta, !B$
\RightLabel{\RightR{\land}}
\BinaryInf$ \Gamma \fCenter \Delta, !A \land !B $
\end{prooftree}
Dalam kasus kita, premis-premisnya harus cocok dengan sekuen terakhir
dari !!{derivation}s yang berakhir pada premis-premis tersebut. Artinya,
$\Gamma$ adalah $!C, !D$, $\Delta$ kosong, $!A$ adalah $!C$, dan $!B$
adalah $!D$. Jadi, agar inferensi itu benar, kesimpulannya ialah $!C, !D
\Sequent !C \land !D$.
\begin{prooftree}
\Axiom$ !C \fCenter !C $
\RightLabel{\LeftR{\Weakening}}
\UnaryInf$ !D, !C \fCenter !C$
\RightLabel{\LeftR{\Exchange}}
\UnaryInf$ !C, !D \fCenter !C$
\Axiom$ !D \fCenter !D $
\RightLabel{\LeftR{\Weakening}}
\UnaryInf$ !C, !D \fCenter !D$
\RightLabel{\RightR{\land}}
\BinaryInf$ !C, !D \fCenter !C \land !D $
\end{prooftree}
Tentu saja, kita juga dapat membalik urutan premis-premisnya; dalam hal
itu $!A$ adalah $!D$ dan $!B$ adalah~$!C$.
\begin{prooftree}
\Axiom$ !D \fCenter !D $
\RightLabel{\LeftR{\Weakening}}
\UnaryInf$ !C, !D \fCenter !D$
\Axiom$ !C \fCenter !C $
\RightLabel{\LeftR{\Weakening}}
\UnaryInf$ !D, !C \fCenter !C$
\RightLabel{\LeftR{\Exchange}}
\UnaryInf$ !C, !D \fCenter !C$
\RightLabel{\RightR{\land}}
\BinaryInf$ !C, !D \fCenter !D \land !C $
\end{prooftree}
\end{ex}

\end{document}
